\scanpage{15}
\mspage{- 104 -}
% archive sheet number: 14
% lines 1–2
$\operatorname{codim}_{Z_1}(Z_{\unsure{o|0}})=1$, il résulte de b) que $\operatorname{codim}_{Z_n}(Z_1)=\operatorname{codim}_{Z_n}(Z_o)-1\handwritten{>}n-1$ ce qui contredit l'hypothèse de récurrence.

% lines 3–4
Lorsque les conditions de la prop.4 sont remplies, on dit que $X$ stisfait à la \uline{condition des chaînes}.

% lines 5–6
\uline{Proposition 5 .- Soit} $X$ \uline{un espace noethé\unsure{r}ien, de dimension finie, vérifiant l'axiome} $(T_o)$. \uline{Les conditions suivantes sont équivalente\handwritten{s}}.

% lines 7–8
a) \uline{Deux chaînes maximales de parties fermées irréductibles de} $X$ \uline{on\handwritten{t} même longueur}.

% lines 9–11
b) \uline{\deleted{\unsure{Toutes les composantes irréductibles} \illegible{4}} \inserted{$X$ est équidimensionnel \deleted{\handwritten{et}} \handwritten{et équicodimensionnel}}, \deleted{les poin\handwritten{ts} fermés de $X$ ont même codimension}, et} $X$ \uline{satisfait à la condition des chaînes}.

% lines 12–13
c) \deleted{\unsure{Les composantes irréductibles de} \illegible{3} \unsure{de} $X$} $X$ \uline{est équidimensionnel et pour deux parties fermées irréductibles} $Y\handwritten{\subset}Z$ \uline{de} $X$, \uline{on a}

% lines 14
\[
 \dim(Z)=\dim(Y)+\operatorname{codim}_Z(Y).\tag*{(4)}
\]

% lines 15–16
d) \inserted{\handwritten{$X$ est équicodimensionnel}} \uline{\deleted{Les codimensions des points fermés de $X$ sont les mêmes}, et pou\handwritten{r} deux parties fermées irréductibles} $Y\handwritten{\subset}Z$ \uline{de} $X$, \uline{on a}

% lines 17
\[
 \operatorname{codim}_X(Y)=\operatorname{codim}_Z(Y)+\operatorname{codim}_X(Z).\tag*{(5)}
\]

% lines left of 17–24
% slanted
\margin{\handwritten{\uline{\unsure{Dire|Dim}} \uuline{Cor}. Soit $X_i$ une comp. irréd. de $X$, et $x$ un pt fermé de $X$. On a $\dim X=\operatorname{codim}_X x=\dim X_i$.}}

% lines 18–25
Les hypothèses sur $X$ entraîne\unsure{n}t que les extrémités d'une chaîne maximale de parties fermées irréductibles sont nécessairement une composante irréductible de $X$ et un point fermé ; en outre, toute chaîne saturée \deleted{\illegible{1}} \inserted{d'extrémités} $Y$ et $Z$ \inserted{($Y\handwritten{\subset}Z$)} est contenue dans une chaîne maximale dont les éléments \inserted{autres que ceux de la chaîne donnée} sont ou bien contenus dans $Y$ ou bien contiennent $Z$. Ces remarques établissent aussitôt l'équivalence de a) et b),et prouvent aussi que \inserted{si a) est vérifiée,} pour toute partie fermée irréductible $Y$ de $X$, o\handwritten{n} a

% lines 26
\[
 \tikz[baseline=(eq.base)]{
  \node[inner xsep=4pt,inner ysep=2pt] (eq) {$\dim(Y)+\operatorname{codim}_X(Y)=\dim(X),$};
  \draw[thin] plot[smooth cycle] coordinates {(eq.west) (eq.north west) (eq.north) (eq.north east) (eq.east) (eq.south east) (eq.south) (eq.south west)};
  \draw[thin] ([xshift=-28mm,yshift=1mm]eq.west) to[bend left=3] (eq.west);
}\tag*{(6)}
\]

% lines 27–29
et comme (3) a aussi lieu, (\deleted{\unsure{5}}\handwritten{4}) et (\deleted{\unsure{6}}\handwritten{5}) s'en déduisent aussitôt. Inversement, (4) entraîne (3), donc la condition des chaînes (prop.\handwritten{4)}) et en outre, en appliquant (4) pour \deleted{\illegible{1}} $Y$ \deleted{\illegible{1}} réduit à un point
